\newcounter{english}
\documentclass{article}

% packages
\usepackage{amsmath, amsthm, thmtools, amsfonts, amssymb, luacode, catchfile, hyperref, ifthen, graphicx, bbm}
\ifcsname c@kobocompile\endcsname
	\usepackage[a5paper, total={1072pt, 1448pt}, margin=10pt, includeheadfoot]{geometry} % set page margins
\else
	\usepackage[a4paper, margin=50pt, includeheadfoot]{geometry}
\fi
\usepackage[shortlabels]{enumitem}
\usepackage[skip=3pt, indent=0pt]{parskip}

% language
\usepackage[bidi=basic, layout=tabular, provide=*]{babel}
\ifcsname c@english\endcsname
	\babelprovide[main, import]{english}
\else
	\babelprovide[main, import, alph=letters.gershayim]{hebrew}
	\setlist[enumerate,2]{label={\alph*.}}
	\babelprovide{rl}
\fi
\babelfont{rm}[Renderer=Harfbuzz]{Libertinus Serif}
\babelfont{sf}{Libertinus Sans}
\babelfont{tt}{Libertinus Mono}

% style
\ifcsname c@stoppagebreaks\endcsname
\else
	\AddToHook{cmd/section/before}{\clearpage}	% Add line break before section
\fi
\linespread{1.3}
\setcounter{secnumdepth}{0}		% Remove default number tags from sections, this won't do well with theorems
\AtBeginDocument{\setlength{\belowdisplayskip}{3pt}}
\AtBeginDocument{\setlength{\abovedisplayskip}{3pt}}
\graphicspath{ {../images/} }

% operators
\DeclareMathOperator\cis{cis}
\DeclareMathOperator\Sp{Span}
\DeclareMathOperator\tr{tr}
\DeclareMathOperator\im{Im}
\DeclareMathOperator\re{Re}
\DeclareMathOperator\diag{diag}
\DeclareMathOperator*\lowlim{\underline{lim}}
\DeclareMathOperator*\uplim{\overline{lim}}
\DeclareMathOperator\rng{rng}
\DeclareMathOperator\Sym{Sym}
\DeclareMathOperator\Arg{Arg}
\DeclareMathOperator\Log{Log}
\DeclareMathOperator\dom{dom}
\DeclareMathOperator\supp{Supp}
\DeclareMathOperator\var{Var}
\DeclareMathOperator\cov{Cov}
\DeclareMathOperator\End{End}
\DeclareMathOperator\id{id}
\DeclareMathOperator\cl{cl}
\DeclareMathOperator\gal{Gal}
\DeclareMathOperator\dist{dist}
\DeclareMathOperator\form{form}
\DeclareMathOperator\term{term}
\DeclareMathOperator\sent{sent}
\DeclareMathOperator\Th{Th}
\DeclareMathOperator\Hom{Hom}
\DeclareMathOperator\Aut{Aut}
\DeclareMathOperator\Char{char}
\DeclareMathOperator\vol{vol}
\DeclareMathOperator\GL{GL}
\DeclareMathOperator\SL{SL}
\DeclareMathOperator\lcm{lcm}

% commands
\newcommand{\cupdot}{\mathbin{\mathaccent\cdot\cup}}
\newcommand{\indicator}{\mathbbm{1}}
\newcommand{\AAA}[0]{\mathbb{A}}
\newcommand{\Aa}[0]{\mathcal{A}}
\newcommand{\BB}[0]{\mathbb{B}}
\newcommand{\Bb}[0]{\mathcal{B}}
\newcommand{\CC}[0]{\mathbb{C}}
\newcommand{\Cc}[0]{\mathcal{C}}
\newcommand{\DD}[0]{\mathbb{D}}
\newcommand{\Dd}[0]{\mathcal{D}}
\newcommand{\EE}[0]{\mathbb{E}}
\newcommand{\FF}[0]{\mathbb{F}}
\newcommand{\Ff}[0]{\mathcal{F}}
\newcommand{\Ii}[0]{\mathcal{I}}
\newcommand{\Gg}[0]{\mathcal{G}}
\newcommand{\HH}[0]{\mathbb{H}}
\newcommand{\Hh}[0]{\mathcal{H}}
\newcommand{\Kk}[0]{\mathcal{K}}
\newcommand{\Ll}[0]{\mathcal{L}}
\newcommand{\Mm}[0]{\mathcal{M}}
\newcommand{\NN}[0]{\mathbb{N}}
\newcommand{\Nn}[0]{\mathcal{N}}
\newcommand{\PP}[0]{\mathbb{P}}
\newcommand{\Pp}[0]{\mathcal{P}}
\newcommand{\QQ}[0]{\mathbb{Q}}
\newcommand{\RR}[0]{\mathbb{R}}
\newcommand{\Rr}[0]{\mathcal{R}}
\newcommand{\Ss}[0]{\mathcal{S}}
\newcommand{\TT}[0]{\mathbb{T}}
\newcommand{\Tt}[0]{\mathcal{T}}
\newcommand{\UU}[0]{\mathbb{U}}
\newcommand{\Uu}[0]{\mathcal{U}}
\newcommand{\Vv}[0]{\mathcal{V}}
\newcommand{\Ww}[0]{\mathcal{W}}
\newcommand{\Xx}[0]{\mathcal{X}}
\newcommand{\Yy}[0]{\mathcal{Y}}
\newcommand{\ZZ}[0]{\mathbb{Z}}
\newcommand{\Zz}[0]{\mathcal{Z}}
\newcommand{\acts}[0]{\circlearrowright}
\newcommand{\explain}[2] {
	\begin{flalign*}
		 && \text{#2} && \text{#1}
	\end{flalign*}
}
\newcommand{\maketitleprint}[1][green]{
	\begin{center}
		\includegraphics[width=6cm]{dragon_#1}
	\end{center}
}

% theorem commands
\newtheoremstyle{c_remark}
	{}	% Space above
	{}	% Space below
	{}% Body font
	{}	% Indent amount
	{\bfseries}	% Theorem head font
	{}	% Punctuation after theorem head
	{.5em}	% Space after theorem head
	{\thmname{#1}\thmnumber{ #2}\thmnote{ \normalfont{\text{(#3)}}}}	% head content
\newtheoremstyle{c_definition}
	{3pt}	% Space above
	{3pt}	% Space below
	{}% Body font
	{}	% Indent amount
	{\bfseries}	% Theorem head font
	{}	% Punctuation after theorem head
	{.5em}	% Space after theorem head
	{\thmname{#1}\thmnumber{ #2}\thmnote{ \normalfont{\text{(#3)}}}}	% head content
\newtheoremstyle{c_plain}
	{3pt}	% Space above
	{3pt}	% Space below
	{\itshape}% Body font
	{}	% Indent amount
	{\bfseries}	% Theorem head font
	{}	% Punctuation after theorem head
	{.5em}	% Space after theorem head
	{\thmname{#1}\thmnumber{ #2}\thmnote{ \text{(#3)}}}	% head content

\ifcsname c@english\endcsname
	\theoremstyle{plain}
	\newtheorem{theorem}{Theorem}[section]
	\newtheorem*{theorem*}{Theorem}
	\newtheorem{lemma}[theorem]{Lemma}
	\newtheorem{proposition}[theorem]{Proposition}
	\newtheorem*{proposition*}{Proposition}

	\theoremstyle{definition}
	\newtheorem{definition}[theorem]{Definition}
	\newtheorem*{definition*}{Definition}
	\newtheorem{example}{Example}[section]
	\newtheorem{exercise}{Exercise}[section]

	\theoremstyle{remark}
	\newtheorem*{remark}{Remark}
	\newtheorem*{solution}{Solution}
	\newtheorem{corollary}[theorem]{Corollary}
	\newtheorem{notation}[theorem]{Notation}
\else
	\theoremstyle{c_plain}
	\newtheorem{theorem}{משפט}[section]
	\newtheorem*{theorem*}{משפט}
	\newtheorem{lemma}[theorem]{למה}
	\newtheorem{proposition}[theorem]{טענה}
	\newtheorem*{proposition*}{טענה}

	\theoremstyle{c_definition}
	\newtheorem{definition}[theorem]{הגדרה}
	\newtheorem*{definition*}{הגדרה}
	\newtheorem{example}{דוגמה}[section]
	\newtheorem{exercise}{תרגיל}[section]

	\theoremstyle{c_remark}
	\newtheorem*{remark}{הערה}
	\newtheorem*{solution}{פתרון}
	\newtheorem{corollary}[theorem]{מסקנה}
	\newtheorem{notation}[theorem]{סימון}
\fi

% Questions related commands
\newcounter{question}
\setcounter{question}{1}
\newcounter{sub_question}
\setcounter{sub_question}{1}

\ifcsname c@english\endcsname
	\newcommand{\question}[1][0]{
		\ifthenelse{#1 = 0}{}{\setcounter{question}{#1}}
		\section{Question \arabic{question}}
		\addtocounter{question}{1}
		\setcounter{sub_question}{1}
	}

	\newcommand{\subquestion}[1][0]{
		\ifthenelse{#1 = 0}{}{\setcounter{sub_question}{#1}}
		\subsection{Part \alph{sub_question}}
		\addtocounter{sub_question}{1}
	}
\else
	\newcommand{\question}[1][0]{
		\ifthenelse{#1 = 0}{}{\setcounter{question}{#1}}
		\section{שאלה \arabic{question}}
		\addtocounter{question}{1}
		\setcounter{sub_question}{1}
	}

	\newcommand{\subquestion}[1][0]{
		\ifthenelse{#1 = 0}{}{\setcounter{sub_question}{#1}}
		\subsection{סעיף \localecounter{letters.gershayim}{sub_question}}
		\addtocounter{sub_question}{1}
	}
\fi

% import lua and start of document
\directlua{common = require ('../common')}

\GetEnv{AUTHOR}

% headers
\author{\AUTHOR}
\date\today

\title{Lie Groups --- Notes}
\setcounter{secnumdepth}{2}
% chktex-file 9
% chktex-file 17

\usepackage{fancyhdr}
\pagestyle{fancy}
\renewcommand{\headrulewidth}{0pt}

\begin{document}
\maketitle
\maketitleprint[teal]{}

\tableofcontents

\section{lecture 1 --- 20.10.2026}
We will follow the book ``Lectures on Lie groups'' by J. Frank Adams through chapter 5.56.

\subsection{Introduction}
Suppose that $V, W$ are finite dimension vector spaces over $\RR$.
Notice that in such spaces all norms are equivalent.
a function $f : U \to W$ where $U \subseteq V$ is differentiable if for all $x \in U$ there exists a linear map $df \mid_x = f'(x) : V \to W$ such that,
\[
	f(x + h)
	= f(x) + (df \mid_x)(h) + o(|h|)
\]
Therefore $f' : U \to \hom(V, W)$ defined by $x \mapsto f'(x)$.
We continue to define in such manner the next differentials recursively, $f'' : U \to \hom(V, \hom(V, W))$ etc.
A function $f$ is $C^\infty$ (smooth) if the derivatives of any order exists and are continuous.
\begin{definition}[Chart]
	Let $X$ be a topological space.
	\textbf{Chart} on $X$ modeled on $V$, $\varphi_{\alpha} : X_{\alpha} \to \Omega_{\alpha}$ is a homeomorphism such that $X_{\alpha} \subseteq X$ is open and $\Omega_{\alpha} \subseteq V$ is open as well.
\end{definition}
In Adams book chart is defined the other way around ($\Omega_{\alpha} \to X_{\alpha}$).
Compatible charts $\varphi_{\alpha}, \varphi_{\beta}$ are charts such that there is a homeomorphism between them and are invertible. \\
An atlas on $X$ is $\{ \varphi_{\alpha} : X_{\alpha} \to \Omega_{\alpha} \}$ a set of charts, such that $\forall \alpha, \beta,\ \varphi_{\beta} \circ \varphi_{\alpha}^{-1}, \varphi_{\alpha} \circ \varphi_{\beta}^{-1}$
are smooth and $\bigcup_{\alpha} X_{\alpha} = X$.
Two atlases are equivalent if their union is atlas.
\begin{exercise}
	This is an equivalent relation.
\end{exercise}
\begin{definition}[Smooth manifold]
	A (smooth) manifold is a Hausdorff 2nd countable topological space $X$ equipped with an equivalence class of atlases.
\end{definition}
Given manifolds $X, Y$ modeled on vector-spaces $V, W$, $X \times Y$ is a manifold modeled on $V \times W$.
$\{ \varphi_{\alpha} : X_{\alpha} \to \Omega_{\alpha} \}$, $\{ \psi_{\beta} : Y_{\beta} \to \Omega_{\beta}' \}$ are atlases,
then $\{ \varphi_{\alpha} \times \psi_{\beta} : X_{\alpha} \times Y_{\beta} \to \Omega_{\alpha} \times \Omega_{\beta'} \}$ is an atlas in $X \times Y$.
Equivalent atlases on $X$ \& on $Y$  yield equivalent atlases on $X \times Y$.
There exists a unique manifold structure on $X \times Y$ such that $X \times Y \to X$ \& $X \times Y \to Y$ are smooth and for any $Z$ manifold and $Z \to X \times Y$ is smooth iff the compositions are smooth.
\begin{definition}[smooth map between manifolds]
	Given $X, Y$ are manifolds modeled on $V, W$ and atlases $\{ \varphi_{\alpha} \}$ and $\{ \psi_{\beta} \}$.
	We say that $f : X \to Y$ is smooth if $\forall \alpha, \beta,\ \varphi_{\beta} \circ f \circ \psi_{\alpha}^{-1}$ is smooth.
\end{definition}
\begin{definition}[Lie group]
	A Lie group $G$ is a set equipped with:
	\begin{enumerate}
		\item a manifold structure
		\item a group structure:
			a group operation $\mu : G \times G \to G$ defined as $\mu(a, b) = a \cdot b$ and inversion $i : G \to G$ defined as $i(x) = x^{-1}$.
	\end{enumerate}
	such that $\mu$ and $i$ are smooth.
\end{definition}
\begin{definition}[Homomorphism of Lie groups]
	Let $G, H$ be Lie groups, a function $f : G \to H$ that is smooth and a group Homomorphism is called Lie group Homomorphism.
\end{definition}
\begin{example}
	$(\RR^n, +)$ is a Lie group.
\end{example}
\begin{example}
	Define $(\RR^n / \ZZ^n, +)$.
	In this case we can define $\forall x \in \RR^n$ and $0 < \varepsilon < \frac{1}{2}$ we can define $\Omega_x = B(x, \varepsilon) \to X_x$ where $X_x$ is its image in $\RR^n / \ZZ^n$.
\end{example}
\begin{example}
	Let $V$ be a vector-space over $\RR$ or $\CC$ denoted as $\Lambda$, and let $\End(V) = \hom_{\Lambda}(V, V)$ ($\End(V) \cong \RR^{n \times n}$),
	and $\aut(V) \subseteq \End(V)$ be the subset of invertible $\Lambda$-linear maps.
	$\aut(V)$ is an open subset of $\End(V)$. In coordinates: $a, b \mapsto a \cdot b$ as matrix multiplication is a polynomial, hence smooth.
	It follows from Kramer's rule that $a \mapsto a^{-1}$ the matrix inversion is achieved by $a^{-1} = \frac{1}{\det a} \operatorname{adj}(a)$, and so it's smooth as well.
\end{example}
\begin{remark}
	For $X, Y$ are manifolds and $f : X \to Y$ is smooth if and only if any $n$ and $S \subseteq \RR^n$ open \& smooth map $S \to X$, the composition $S \to X \to Y$ is smooth,
	if and only if for all $g : Y \to \RR$ smooth, $f \circ g$ is smooth.
\end{remark}
\begin{definition}[Tangent space]
	Tangent space to a manifold $X$ at $x \in X$ is defined as,
	\[
		T_x X
		= \{ \gamma : (-\delta, \delta) \to X \text{ smooth}, \gamma(0) = x \} / \sim
	\]
	where $\gamma_1 \sim \gamma_2 \iff \forall f : U \to \RR$ where $x \in U$ is open neighborhood, smooth map, then $\frac{d}{dt} \mid_{t = 0} f \circ \gamma_1 = \frac{d}{dt} \mid_{t = 0} f \circ \gamma_2$.
\end{definition}
\begin{example}
	For $x \in X \subseteq V$ open, namely an open subset of a vector space, $T_x X = V$ via $\gamma(t) = x + t v$ for $v \in V$.
\end{example}
Smooth map of manifolds $f : X \to Y$ induces $f' = f_* = df \mid_x : T_x X \to T_{f(x)} Y$. \\
Given a manifold $X$ \& atlas $\{ \varphi_{\alpha} : X_{\alpha} \to \Omega_{\alpha} \}$,
\[
	{(\varphi_{\alpha})}^* : T_x X_{\alpha} \to T_{\varphi_{\alpha}(x)} \Omega_{\alpha} = V
\]
Given a well defined vector-spaces structure on $T_x X$.
\begin{definition}[Lie algebra]
	Suppose that $G$ is a Lie group and $e \in G$ its unit, we define $L(G) = T_e G$ (as a vector-spaces),
	Lie algebra of $G$.
	Also notated as $\operatorname{Lie}(G)$ or $\mathfrak{G}$.
\end{definition}
Let $G$ be a Lie group, $x \in G$, and define left multiplication,
\[
	L_x : G \to x g
\]
defined a diffeomorphism $G \to G$ by $e \mapsto x$ hence a linear isomorphism $T_e G \xrightarrow{\sim} T_x G$,
then left trivial injection $TG \cong G \times L(G)$.

\listoftheorems[title=definitions \& theorems,ignoreall,show={theorem,definition},swapnumber,onlynamed={proposition}]
\end{document}
